\begin{resumen}
	El presente trabajo de diploma presenta la propuesta de aplicación web DermaUH 3.0, 
	desarrollada con el objetivo de suplir diversas dificultades de la versión anterior e 
	incluir mejoras tanto de rendimiento como en sus modelos de clasificación de imágenes.

	Para el proceso de investigación se analizaron las diferentes herramientas existentes 
	en el mercado que utilizan inteligencia artificial para los diagnósticos asistidos, 
	atendiendo a sus ventajas y desventajas para los doctores cubanos, y se tomaron ideas 
	para el desarrollo de la interfaz de usuario. Posteriormente se revisaron diferentes 
	sitios académicos especializados en la colección de imágenes dermatoscópicas, 
	argumentando por qué estos no podían ser utilizados para entrenar los nuevos modelos 
	de aprendizaje de máquinas. Finalmente se realizó un estudio de diferentes usos de los 
	modelos explicativos de inteligencia artificial (XAI) para obtener experiencia sobre 
	cómo integrarlos a una aplicación web.

	A raíz de la investigación descrita anteriormente, se llegó a una propuesta de 
	solución que cubre las necesidades planteadas para la nueva versión de DermaUH 
	haciendo uso de Django, React y PostgreSQL como tecnologías principales. Se siguió el 
	patrón de Arquitectura por Capas para el desarrollo de la aplicación, donde cada nivel 
	presenta su propia propuesta de diseño. Para el frontend se utilizó una arquitectura 
	basada en características y para el backend se implementaron varias aplicaciones 
	Django almacenadas en un mismo servidor.

	El entrenamiento de los nuevos modelos se realizó con un conjunto de imágenes de 
	pieles cubanas proporcionado por las autoridades del ministerio de Salud Pública de 
	Cuba (MINSAP). Esta acción fue llevada a cabo con la herramienta Kaggle para ganar 
	en tiempo de ejecución debido al potente procesador que brinda dicha herramienta.

	Finalizado el proceso de desarrollo, el principal resultado fue la obtención de una 
	aplicación web capaz de brindar diagnóstico asistido por Inteligencia Artificial. 
	El modelo utilizado fue entrenado con pieles cubanas y se vinculó con un modelo 
	explicativo para justificar los diagnósticos brindados.
\end{resumen}

\begin{abstract}
	This diploma thesis presents the proposed web application DermaUH 3.0, developed to 
	address various shortcomings of the previous version and include improvements in 
	both performance and image classification models.

	For the research process, the different existing tools
	on the market that use artificial intelligence for computer-assisted diagnoses 
	were analyzed, considering their advantages and disadvantages for Cuban physicians, 
	and ideas were gathered for the development of the user interface. Subsequently, 
	various academic websites specializing in the collection of dermatoscopic 
	images were reviewed, and arguments were presented as to why these could not be used 
	to train the new machine learning models. Finally, a study of different uses of
	explanatory models of artificial intelligence (XAI) was conducted to gain experience
	on how to integrate them into a web application.

	As a result of the research described above, a proposed solution was developed 
	that addresses the needs identified for the new version of DermaUH, using Django, 
	React, and PostgreSQL as the main technologies. The Layered Architecture pattern was 
	followed for the application's development, where each layer presents its own design 
	proposal. A feature-based architecture was used for the frontend, and several Django 
	applications stored on the same server were implemented for the backend.

	The training of the new models was performed using a set of Cuban skin images provided
	by the Cuban Ministry of Public Health (MINSAP). This was done using the Kaggle tool 
	to improve execution time due to its powerful processor capabilities.

	Upon completion of the development process, the main result was a web application 
	capable of providing AI-assisted diagnosis. The model used was trained with Cuban 
	skins and was linked to an explanatory model to justify the diagnoses provided.
\end{abstract}